\medskip

\emph{Principaux domaines abordés}: Étude des fonctions.
Fonction logarithme.

\bigskip

\textbf{Partie A}

\begin{center}
\psset{unit=0.6cm,arrowsize=2pt 3}
\begin{pspicture*}(-4,-5)(15.5,8)
\psgrid[gridlabels=0pt,subgriddiv=1,gridwidth=0.15pt]
\psaxes[linewidth=1.25pt,labelFontSize=\scriptstyle]{->}(0,0)(-4,-5)(15.5,8)
\psplot[plotpoints=2000,linewidth=1.25pt,linecolor=blue]{3.001}{15.5}{x dup mul x sub 6 sub ln}
\psplot[plotpoints=2000,linewidth=1.25pt,linecolor=red,linestyle=dashed]{3.001}{15.5}{x 2 mul 1 sub x dup mul x sub 6 sub div}
\uput[r](3.1,7.7){\red $\mathcal{C}_2$}\uput[r](3.1,-4){\blue $\mathcal{C}_1$}
\end{pspicture*}
\end{center}

Dans le repère orthonormé ci-dessus, sont tracées les courbes représentatives d'une fonction $f$ et de sa fonction dérivée, notée $f'$, toutes deux définies sur $]3~;~+\infty[$.

\medskip

\begin{enumerate}
\item Associer à chaque courbe la fonction qu'elle représente. Justifier.
\item Déterminer graphiquement la ou les solutions éventuelles de l'équation $f(x) = 3$. 
\item Indiquer, par lecture graphique, la convexité de la fonction $f$.
\end{enumerate}

\bigskip

\textbf{Partie B}

\medskip

\begin{enumerate}
\item Justifier que la quantité $\ln \left(x^2- x- 6\right)$ est bien définie pour les valeurs $x$ de l'intervalle $]3~;~ +\infty[$, que l'on nommera $I$ dans la suite.
\item On admet que la fonction $f$ de la Partie A est définie par $f(x) = \ln \left(x^2- x- 6\right)$ sur $I$. 

Calculer les limites de la fonction $f$ aux deux bornes de l'intervalle $I$.

En déduire une équation d'une asymptote à la courbe représentative de la fonction $f$ sur $I$.
\item 
	\begin{enumerate}
		\item Calculer $f'(x)$ pour tout $x$ appartenant à $I$.
		\item Étudier le sens de variation de la fonction $f$ sur $I$.
		
Dresser le tableau des variations de la fonction $f$ en y faisant figurer les limites aux bornes de $I$.
	\end{enumerate}
\item 
	\begin{enumerate}
		\item Justifier que l'équation $f(x) = 3$ admet une unique solution $\alpha$ sur l'intervalle ]5~;~ 6[.
		\item Déterminer, à l'aide de la calculatrice, un encadrement de $\alpha$ à $10^{-2}$ près.
	\end{enumerate}
\item 
	\begin{enumerate}
		\item Justifier que $f''(x) = \dfrac{- 2x^2 + 2x - 13}{\left(x^2 - x - 6\right)^2}$.
		\item Étudier la convexité de la fonction $f$ sur $I$.
	\end{enumerate}
\end{enumerate}

\bigskip
